% FORMALIZACION_REUNIDA: alineamiento dodecafasico y lector angular.
% CERTIFICADO_NUEVO: composicion exacta conservada en el suplemento local.
% La generacion previa del sello y la seleccion sectorial mantienen sus pruebas.
\section{Réalisation équivariante du lecteur angulaire dans l’incidence de Witt}
\label{sec:ii-enlace-angular-witt}

La coordinatisation angulaire et l’incidence exceptionnelle conservent des informations différentes
du même développement APP--TRIT--TPK. La première exprime des coordonnées de
rotation ; la seconde conserve des relations entre positions, supports et
orientations. Leur comparaison nécessite un transport qui préserve ces
relations et permette de retrouver les coordonnées complètes. Nous construisons ici
ce transport et démontrons sa compatibilité avec les permutations
orientées du triplet régional.

Les antécédents de la composition sont le sceau dodécaphasique \(K\), sa
face marquée \(B_K\), le secteur du projecteur isotypique tridimensionnel
\(P_3^{\mathrm{ico}}\), les branches
compatibles de la filtration nonadique et le système sectoriel de la
Définition~\ref{def:ii09-sistema-sectorial}. Chacun intervient avec sa
structure, et non uniquement par sa valeur scalaire. Les espaces
vectoriels réels de cette partie sont des réalisations ultérieures de ces
secteurs déjà exprimés. La génération du sceau, le transport de ses
coordonnées et l’évaluation du lecteur correspondent à des étapes distinctes de
la composition causale.

\subsection{Triplet régional et droite d’orientation}

Soit \(V\) l’espace de comptage des trois régions positives de
\(\pi_{\rm HMT}\), de base orthonormée ordonnée
\[
 (v_1,v_2,v_3)=(933555,393555,339555).
\]
La notation identifie chaque vecteur de base à la région qu’il représente.
Dans la fenêtre de prolongement d’indice sept, les onze branches admises
se répartissent en orbites de tailles \(3+3+1+1+3\) sous la permutation
simultanée des trois positions finales des deux premiers mots.
Le secteur qui conserve une continuation à l’horizon deux se compose de
\[
\begin{array}{c|c|c}
\text{branche}&\text{triplet de mots}&\text{région positive}\\ \hline
0&200112\mid212221\mid110110&339555\\
1&200121\mid212212\mid110110&393555\\
2&200211\mid212122\mid110110&933555.
\end{array}
\]
La branche centrale est celle qui se prolonge également à l’horizon trois. Ces
horizons appartiennent à la même filtration de prolongements compatibles ;
les cardinalités successives du support sont \(11\), \(3\) et \(1\).

L’application de la queue du premier mot à la région positive
substitue coordonnée par coordonnée \(1\mapsto3\) et \(2\mapsto9\), et ajoute
la queue régionale \(555\). Par conséquent, elle conserve la position du symbole
exceptionnel et commute avec toute permutation des trois positions. L’ordre
des branches est \((339555,393555,933555)\) ; le passage depuis l’ordre de
\(V\) est la permutation \(R=(1\,3)\).

La droite d’orientation de \(V\) est \(o(V)=\det V\). L’application
\begin{equation}
 \Psi:\Lambda^2V\longrightarrow\operatorname{Hom}(V,\det V),
 \qquad \Psi(u\wedge v)(w)=u\wedge v\wedge w
 \label{eq:ii-witt-hodge-natural}
\end{equation}
s’identifie au moyen de la métrique de comptage à une isométrie
\(\Lambda^2V\to o(V)\otimes V\). Pour
\(\omega=v_1\wedge v_2\wedge v_3\), ses valeurs sur la base extérieure
sont
\[
 v_1\wedge v_2\mapsto\omega\otimes v_3,
 \qquad v_2\wedge v_3\mapsto\omega\otimes v_1,
 \qquad v_1\wedge v_3\mapsto-\omega\otimes v_2.
\]
Dans l’ordre angulaire \((12,23,13)\), le transport conjoint de la base et de
l’orientation a pour matrice
\begin{equation}
 D=\det(R)P_R
 \begin{pmatrix}0&1&0\\0&0&-1\\1&0&0\end{pmatrix}
 =\begin{pmatrix}-1&0&0\\0&0&1\\0&-1&0\end{pmatrix}.
 \label{eq:ii-witt-orientation-D}
\end{equation}
Ainsi,
\(D(\theta_{12},\theta_{23},\theta_{13})^{\mathsf T}
=(-\theta_{12},\theta_{13},-\theta_{23})^{\mathsf T}\).
Dans toute cette partie, \(P_p e_j=e_{p(j)}\) définit la convention des
matrices de permutation, et l’astérisque désigne l’adjoint pour les métriques
de comptage déclarées.

\begin{proposition}[Naturalité de l’identification orientée]
\label{prop:ii-witt-D-natural}
Pour toute \(p\in S_3\),
\[
 D\Lambda^2P_p=\det(p)P_{RpR^{-1}}D,
 \qquad D^*D=I_3.
\]
\end{proposition}
\begin{proof}
L’identité
\((P_pu)\wedge(P_pv)\wedge(P_pw)=\det(p)(u\wedge v\wedge w)\)
prouve la naturalité de \eqref{eq:ii-witt-hodge-natural}. Le changement
simultané de l’ordre régional introduit la conjugaison par \(R\).
Enfin, les colonnes de \(D\) sont des vecteurs de coordonnées signés,
d’où \(D^*D=I_3\).
\end{proof}

La droite d’orientation explique la différence entre le caractère
\((3,1,0)\) du triplet de positions et le caractère \((3,-1,0)\)
du triplet orienté, dans les classes de l’identité, des transpositions et des cycles
d’ordre trois. Les trois transpositions fournissent un contrôle explicite :
lorsqu’on omet le facteur déterminant, l’identité précédente change de signe.
Un changement global de l’ancre orientée se transporte, en revanche, comme changement
conjoint de coordonnées.

\subsection{Secteur dodécaphasique et sélection variationnelle de l’orientation}

La coordinatisation du sceau précédemment exprimé est
\begin{equation}
 K=(234,543,140,729,659,824,621,58,914,794,146,601).
 \label{eq:ii-witt-K-marked}
\end{equation}
Pour expliciter sa réalisation d’incidence, nous écrivons la matrice ternaire
de Paley dans la jauge utilisée :
\[
 A_W=\begin{pmatrix}
 0&1&1&1&1&1\\
 1&0&1&2&2&1\\
 1&1&0&1&2&2\\
 1&2&1&0&1&2\\
 1&2&2&1&0&1\\
 1&1&2&2&1&0
 \end{pmatrix},
 \qquad C(a)=(a,aA_W),\quad a\in\mathbb F_3^6.
\]
Le support est pris sur les douze coordonnées de \(C(a)\).
L’intersection marquée a les deux expressions compatibles
\begin{equation}
 B_K=\operatorname{supp}C(010211)\cap\operatorname{supp}C(201101)
     =\{p:K_p\ge729\}=\{4,6,9,10\}.
 \label{eq:ii-witt-B-marked}
\end{equation}
L’égalité s’évalue au moyen des multiplications ternaires de la matrice
présentée. En particulier, l’intersection s’obtient comme ensemble de
positions, avant d’utiliser son cardinal quatre.

Pour fixer complètement la représentation des douze positions,
nous considérons \(A_5\) agissant par multiplication à gauche sur les
classes \(r_pC_5\), où \(C_5=\langle(1\,2\,3\,4\,5)\rangle\).
Les représentants, écrits comme listes d’images sur
\(\{1,\ldots,5\}\), sont
\[
\begin{array}{c|c@{\qquad}c|c}
p&r_p&p&r_p\\ \hline
1&(4,3,5,2,1)&7&(5,4,3,2,1)\\
2&(4,2,1,3,5)&8&(1,3,4,2,5)\\
3&(1,2,4,5,3)&9&(4,2,3,5,1)\\
4&(2,5,3,4,1)&10&(3,2,5,4,1)\\
5&(4,3,1,5,2)&11&(4,5,2,3,1)\\
6&(5,1,2,3,4)&12&(5,3,2,4,1).
\end{array}
\]
Nous notons cette représentation de permutation \(\rho\) et posons
\(V_{11}=\mathbf1^\perp\subset\mathbb R^{12}\). Le projecteur
isotypique tridimensionnel icosaédrique est déterminé par
\begin{equation}
 P_3^{\mathrm{ico}}=\frac3{60}\sum_{a\in A_5}\chi_3(a^{-1})\rho(a),
 \label{eq:ii-witt-projector-character}
\end{equation}
où les valeurs de \(\chi_3\) sur les classes de types
\(1,2^2,3,5_a,5_b\) sont
\[
 \left(3,-1,0,\frac{1+\sqrt5}{2},\frac{1-\sqrt5}{2}\right).
\]
La classe \(5_a\) contient \((1\,2\,3\,4\,5)\), et la classe
\(5_b\) contient son carré. Cette convention distingue les deux
facteurs directs tridimensionnels. L’exposant \(\mathrm{ico}\) identifie
cette composante de la représentation de \(A_5\) ; le projecteur cyclique
\(P_3\) qui intervient dans~\eqref{eq:det} agit sur le transport de douze
positions avec sa définition propre. L’évaluation finie de
\eqref{eq:ii-witt-projector-character} satisfait
\(P_3^{\mathrm{ico}}=(P_3^{\mathrm{ico}})^*=(P_3^{\mathrm{ico}})^2\), \(P_3^{\mathrm{ico}}\mathbf1=0\) et
\(\operatorname{Tr}P_3^{\mathrm{ico}}=3\).

Les deux vecteurs marqués du secteur sont
\[
 b=P_3^{\mathrm{ico}}\left(\mathbf1_{B_K}-\frac13\mathbf1\right),
 \qquad u=P_3^{\mathrm{ico}}(K-\overline K\mathbf1),
 \qquad \overline K=\frac1{12}\sum_{p=1}^{12}K_p,
 \qquad \|b\|^2=1.
\]
Pour chacun des dix sous-groupes d’ordre trois \(C_3\subset A_5\),
soit \(H=N_{A_5}(C_3)\cong S_3\). Le caractère
\(\operatorname{sgn}_H\) prend la valeur \(+1\) sur \(C_3\) et
\(-1\) sur \(H\setminus C_3\). C’est le caractère du quotient
\(H/C_3\), tandis que le signe ambiant des éléments de \(A_5\)
est toujours positif. On définit
\begin{equation}
 e_{{\rm sgn},H}=\frac16\sum_{a\in H}\operatorname{sgn}_H(a)\rho(a),
 \qquad \mathcal E_B(H)=\|e_{{\rm sgn},H}b\|^2.
 \label{eq:ii-witt-energy}
\end{equation}

\begin{proposition}[Normalisateur et orientation déterminés par le sceau]
\label{prop:ii-witt-variational-selection}
La règle du maximum de \eqref{eq:ii-witt-energy} sélectionne un unique
normalisateur \(H_*\). Son énergie et la deuxième énergie distincte sont
\[
 \frac23+\frac{2\sqrt5}{15},\qquad
 \frac13+\frac{2\sqrt5}{15},
\]
respectivement ; la séparation est \(1/3\). La projection de \(u\)
sélectionne le même normalisateur. Dans \(H_*\), la maximisation de
\(s(a)=\langle b,\rho(a)u\rangle\), séparément parmi les éléments
d’ordre trois et les involutions, détermine
\[
 c_*=(1,2,4,5,3),\qquad r_*=(2,1,5,4,3)
\]
en notation d’images sur cinq points. On a
\(r_*c_*r_*=c_*^{-1}\).
\end{proposition}
\begin{proof}
L’évaluation des sommes finies précédentes dans \(\mathbb Q(\sqrt5)\)
donne la distribution suivante des dix énergies :
\[
\begin{array}{c|c@{\qquad}c|c}
\text{énergie}&\text{multiplicité}&\text{énergie}&\text{multiplicité}\\ \hline
2/3+2\sqrt5/15&1&2/3-2\sqrt5/15&1\\
1/3+2\sqrt5/15&2&1/3-2\sqrt5/15&2\\
1/6+\sqrt5/30&2&1/6-\sqrt5/30&2.
\end{array}
\]
La plus grande appartient au sous-groupe engendré par \(c_*\). La suivante
s’obtient en soustrayant \(1/3\). La même évaluation avec \(u\) a un maximum
unique sur ce sous-groupe, de valeur
\(2526869/12+2813609\sqrt5/30\).
Le score des deux orientations du cycle est
\[
 s(c_*)=\frac{1341}{4}+\frac{2409\sqrt5}{20},\qquad
 s(c_*^{-1})=\frac{357}{2}+\frac{1853\sqrt5}{10}.
\]
Leur différence est \(627/4-1297\sqrt5/20>0\), comme on le vérifie
en comparant les carrés positifs. Les scores des trois
involutions sont
\[
\begin{array}{c|c}
\text{liste des images}&s(a)\\ \hline
(2,1,3,5,4)&-1341/4-2409\sqrt5/20\\
(2,1,4,3,5)&-1659/4-2667\sqrt5/20\\
(2,1,5,4,3)&-357/2-1853\sqrt5/10.
\end{array}
\]
Le maximum unique correspond à \(r_*\). La composition des deux listes
d’images vérifie la relation diédrale. Toutes ces opérations utilisent
\(K\), \(B_K\), \(P_3^{\mathrm{ico}}\) et les actions finies déjà spécifiées.
\end{proof}

Nous écrivons \(\vartheta:S_3\to H_*\) pour l’isomorphisme déterminé
par \(\vartheta((1\,2\,3))=c_*\) et
\(\vartheta((1\,3))=r_*\). La réflexion des positions fixe la branche
centrale qui atteint l’horizon trois. Cette orientation complète
l’identification régionale avant toute évaluation d’angles.

\subsection{Entrelaceur de Reynolds et normalisation polaire}

Dans la base des branches, soit \(\rho_o(g)=\det(g)P_g\) et
\(w=(0,1,0)^{\mathsf T}\). L’opérateur de Reynolds défini par sommation
sur le groupe est
\begin{equation}
 A_R=\sum_{g\in S_3}\rho_o(g)\,|w\rangle\langle b|\,
                       \rho(\vartheta(g))^{-1}.
 \label{eq:ii-witt-reynolds}
\end{equation}
Le symbole \(A_R\) distingue cet opérateur de l’angle
\(A_{\rm deg}\). Chaque terme a une échelle unitaire ; la normalisation
métrique s’effectue au moyen de son facteur polaire.

\begin{proposition}[Isométrie polaire du secteur orienté]
\label{prop:ii-witt-polar}
On a
\[
 A_R=A_RP_3^{\mathrm{ico}},\qquad
 \rho_o(g)A_R=A_R\rho(\vartheta(g)),\qquad
 G=A_RA_R^*>0,\qquad A_R^*G^{-1}A_R=P_3^{\mathrm{ico}}.
\]
Avec \(E=\mathbf1\mathbf1^{\mathsf T}/3\),
\begin{equation}
 G=\lambda_sE+\lambda_t(I_3-E),\qquad
 \lambda_s=8+\frac{8\sqrt5}{5},\quad
 \lambda_t=\frac12-\frac{\sqrt5}{5}>0.
 \label{eq:ii-witt-gram-spectrum}
\end{equation}
En conséquence,
\[
 U_{GK}=G^{-1/2}A_R:P_3^{\mathrm{ico}}V_{11}\longrightarrow o(V)\otimes V
\]
est une isométrie surjective :
\(U_{GK}^*U_{GK}=P_3^{\mathrm{ico}}\) et \(U_{GK}U_{GK}^*=I_3\).
\end{proposition}
\begin{proof}
Le projecteur \(P_3^{\mathrm{ico}}\) commute avec \(\rho\), et \(b=P_3^{\mathrm{ico}}b\) ;
donc \(A_R=A_RP_3^{\mathrm{ico}}\). Réindexer la somme de Reynolds prouve
l’équivariance. La multiplication de la somme finie donne les entrées diagonales
\(3+2\sqrt5/5\) et les entrées non diagonales
\(5/2+3\sqrt5/5\) dans \(G\). Les projections \(E\) et \(I_3-E\)
fournissent \eqref{eq:ii-witt-gram-spectrum}. Les deux valeurs propres sont
positives ; pour la plus petite, \(25>20\) suffit. Par conséquent, \(A_R\) est de
rang trois. Son espace des lignes est contenu dans le secteur de rang trois
de \(P_3^{\mathrm{ico}}\), et coïncide avec lui. Le projecteur orthogonal sur cet espace
des lignes est \(A_R^*(A_RA_R^*)^{-1}A_R\), ce qui donne la dernière identité.
Les deux équations du facteur polaire s’ensuivent par substitution.
\end{proof}

La racine positive conserve l’équivariance parce que les deux projecteurs
spectraux commutent avec la représentation orientée. L’unicité du
facteur polaire correspond à l’opérateur de Reynolds et à l’orientation
déterminés par les règles précédentes.

\subsection{Incidence point--hexade et changement conjoint de jauge}

Soit \(\mathcal W\) le plan \(S(5,6,12)\) de 132 hexades obtenu
à partir des supports de poids six du code ternaire \(C(a)\). La matrice
d’incidence est
\[
 I_{H,p}=\mathbf1_{p\in H},\qquad
 (Iz)_H=\sum_{p\in H}z_p.
\]
Les blocs se reconstruisent en énumérant \(a\in\mathbb F_3^6\),
en retenant le poids six et en réunissant les supports égaux. La vérification
finie que chaque sous-ensemble de cinq points appartient à une unique
hexade fixe la propriété de plan utilisée ci-après.

\begin{lemma}[Identités métrique et spectrale de l’incidence]
\label{lem:ii-witt-incidence-identities}
Soient \(J_{12}\) et \(J_{132,12}\) les matrices de uns indiquées par
leurs dimensions, et soit
\[
 (G_{\mathcal W})_{H,H'}=\binom{|H\cap H'|}{4}.
\]
Alors
\begin{equation}
 I^*I=36I_{12}+30J_{12},\qquad
 G_{\mathcal W}I=30I+15J_{132,12}.
 \label{eq:ii-witt-incidence-identities}
\end{equation}
En particulier, \(U_W=I/6\) est une isométrie sur
\(V_{11}=\mathbf1^\perp\), dont l’image est contenue dans
\(E_{30}=\ker(G_{\mathcal W}-30I_{132})\).
\end{lemma}
\begin{proof}
La propriété de plan donne 66 hexades par point et 30 par paire de points :
\[
 \frac{\binom{11}{4}}{\binom54}=66,
 \qquad \frac{\binom{10}{3}}{\binom43}=30.
\]
Ce sont les entrées diagonales et non diagonales de \(I^*I\).
Pour la seconde identité, on compte les paires \((T,H')\), où
\(T\subset H\cap H'\), \(|T|=4\) et \(p\in H'\).
Si \(p\in H\), dix tétrades \(T\subset H\) contiennent \(p\),
et chacune appartient à quatre hexades ; les cinq autres tétrades,
unies à \(p\), déterminent chacune une hexade. Le résultat est
\(10\cdot4+5=45\). Si \(p\notin H\), chacune des quinze
tétrades détermine avec \(p\) une hexade, et le résultat est 15.
On obtient ainsi \(30I_{H,p}+15\). Les matrices de uns s’annulent sur
\(V_{11}\), ce qui prouve l’isométrie et l’appartenance spectrale.
\end{proof}

Le changement de jauge des douze positions est
\begin{equation}
 \sigma=(1,5,2,6,7,4,3,10,11,9,8,12),
 \label{eq:ii-witt-sigma}
\end{equation}
écrit comme liste d’images avec des indices commençant à un. Si
\(\rho_{\rm perm}\) est l’action ponctuelle sous-jacente à \(\rho\), les
six permutations
\[
 \tau_g=\sigma\rho_{\rm perm}(\vartheta(g))\sigma^{-1},\qquad g\in S_3,
\]
préservent \(\mathcal W\). Leurs deux générateurs ont pour listes d’images
\[
\begin{aligned}
 \tau_{(123)}&=(3,6,5,2,1,4,9,12,11,8,7,10),\\
 \tau_{(13)}&=(2,1,4,3,6,5,8,7,10,9,12,11).
\end{aligned}
\]
L’affirmation se vérifie sur le plan fini : on applique ces deux
permutations génératrices aux supports
\(C(a)\) de poids six et on vérifie que leurs images parcourent les
mêmes 132 supports. Les produits de ces générateurs donnent les quatre
autres permutations. L’action induite \(Q_{\tau_g}\) sur les hexades
satisfait
\begin{equation}
 I P_{\tau_g}=Q_{\tau_g}I,
 \qquad
 P_\sigma\rho(\vartheta(g))=P_{\tau_g}P_\sigma.
 \label{eq:ii-witt-recalibration}
\end{equation}
La première identité exprime exactement l’équivalence
\(p\in H\Longleftrightarrow\tau_g(p)\in\tau_g(H)\).

Ce changement transporte conjointement les positions, le code, l’incidence,
l’origine, la chambre et l’orientation du sceau. Dans la jauge initiale, cinq
éléments de ce sous-groupe d’ordre six ne préservent pas le plan fixe.
La conjugaison \eqref{eq:ii-witt-sigma} fournit la jauge compatible
et maintient ce contrôle sur l’étiquetage initial. La procédure
présente un changement de jauge ; son unicité n’intervient pas dans la preuve.

\subsection{Composition isométrique et récupération des coordonnées}

\begin{theorem}[Transport angulaire dodécaphasique--Witt]
\label{thm:ii-witt-angular-transport}
La composition
\begin{equation}
 \boxed{F=\frac16 I P_\sigma A_R^*G^{-1/2}D:
        \mathbb R^3_{\rm ang}\longrightarrow E_{30}}
 \label{eq:ii-witt-F}
\end{equation}
est une isométrie. Son action en coordonnées et son équivariance sont
\begin{align}
 (F\boldsymbol\theta)_H
 &=\frac16\sum_{p\in H}
       (P_\sigma A_R^*G^{-1/2}D\boldsymbol\theta)_p,
       \label{eq:ii-witt-F-coordinates}\\
 F X_g&=Q_{\tau_g}F,
 \qquad X_g=\Lambda^2P_{R^{-1}gR}.
       \label{eq:ii-witt-F-equivariance}
\end{align}
\end{theorem}
\begin{proof}
Les colonnes de \(A_R^*\) appartiennent à \(P_3^{\mathrm{ico}}V_{11}\), et
\(P_\sigma\) conserve \(V_{11}\). Le
Lemme~\ref{lem:ii-witt-incidence-identities} démontre donc
l’appartenance de l’image de \(F\) à \(E_{30}\). De plus,
\begin{align*}
 F^*F
 &=D^*G^{-1/2}A_RP_\sigma^*
       \frac{I^*I}{36}P_\sigma A_R^*G^{-1/2}D\\
 &=D^*G^{-1/2}A_RA_R^*G^{-1/2}D
  =D^*D=I_3.
\end{align*}
La formule en coordonnées développe la multiplication par l’incidence.
Pour l’équivariance, on compose, dans leur ordre, la naturalité de \(D\),
l’équivariance de \(A_R^*\), la commutation de \(G^{-1/2}\) avec
\(\rho_o\) et les deux identités de
\eqref{eq:ii-witt-recalibration}. Cela produit
\eqref{eq:ii-witt-F-equivariance} pour chaque \(g\in S_3\).
\end{proof}

\begin{proposition}[Covariance de la jauge complète]
\label{prop:ii-witt-gauge-covariance}
Soient \(S,T,V_o\) des changements orthogonaux de coordonnées dans l’espace
dodécaphasique, l’espace ponctuel de Witt et le secteur orienté,
respectivement. Soit \(Q_T\) le changement correspondant de coordonnées
dans l’espace d’incidence ; pour un réétiquetage du plan, c’est la
permutation induite des hexades. Si l’on transporte conjointement
\[
\begin{gathered}
 A_R'=V_oA_RS^*,\quad G'=V_oGV_o^*,\quad
 P_\sigma'=TP_\sigma S^*,\\
 D'=V_oD,\qquad U_W'=Q_TU_WT^*,
\end{gathered}
\]
la composition satisfait \(F'=Q_TF\).
\end{proposition}
\begin{proof}
L’unicité de la racine positive donne
\((G')^{-1/2}=V_oG^{-1/2}V_o^*\). En substituant les cinq
transformations dans
\(F'=U_W'P_\sigma'(A_R')^*(G')^{-1/2}D'\), les facteurs
\(T^*T\), \(S^*S\) et \(V_o^*V_o\) s’annulent par paires, et il reste
\(Q_TU_WP_\sigma A_R^*G^{-1/2}D=Q_TF\).
\end{proof}

La covariance précédente décrit un changement de coordonnées de toutes les
données marquées. Son application à une transformation cardinale active nécessite
l’identité correspondante entre les actions sur l’état ; les deux
notions conservent ainsi leurs domaines propres.

\subsection{Réalisation fidèle du système sectoriel angulaire}

On utilise la coordinatisation \(x_\theta\) de
\eqref{eq:ii-angular-chart},
\[
 x_\theta=\begin{pmatrix}A_{\rm deg}\\C^*_{\rm deg}/6\\
                              (180/\pi_{\rm HMT})\Delta_4\end{pmatrix}.
\]
La composante \(h_\theta=C^*_{\rm deg}/6\) répartit la composante
impaire du registre entre six paires antipodales. La
Définition~\ref{def:ii09-sistema-sectorial} fixe les trois règles
\[
 R_{12}=(2,-p,bq),\quad R_{23}=(0,q,-p^2/2),\quad
 R_{13}=(1,-pb,-(o_8-h_6)/t)
\]
pour \((b,p,q,h_6,o_8,t)=(4,5,7,6,8,3)\). La face de cardinal
quatre, l’ordre pentadique, la coordonnée du pivot, la différence
hexadique--octadique et l’alphabet tritique conservent des types distincts.
La normalisation orbitale bidirectionnelle de
\ref{prop:ii09-semicuadrado-orbital} conserve le terme \(p^2/2\),
y compris la contribution des paires fixes.

Son évaluation, déjà démontrée dans le
Théorème~\ref{thm:ii09-matriz-sectorial}, donne
\[
 \mathsf M_{\rm CKM}=
 \begin{pmatrix}2&-5&28\\0&7&-25/2\\1&-20&-2/3\end{pmatrix},
 \qquad\det\mathsf M_{\rm CKM}=-\frac{3857}{6}.
\]
Le transport \(F\) agit sur les coordonnées exprimées par ces
règles ; le choix antérieur des fonctionnelles conserve sa place dans la
construction sectorielle.

\begin{corollary}[Lecteur angulaire dans le module d’incidence]
\label{cor:ii-witt-composed-reader}
L'application
\begin{equation}
 \mathcal C=F\mathsf M_{\rm CKM}
 \label{eq:ii-witt-composed-reader}
\end{equation}
satisfait
\[
 \mathcal C^*\mathcal C=\mathsf M_{\rm CKM}^*\mathsf M_{\rm CKM},
 \qquad
 \mathsf M_{\rm CKM}^{-1}F^*\mathcal C=I_3.
\]
C’est une isométrie de l’espace de \(x_\theta\) muni de la métrique
\(\mathsf M_{\rm CKM}^*\mathsf M_{\rm CKM}\) sur son image.
\end{corollary}
\begin{proof}
On substitue \(F^*F=I_3\) dans les deux produits. Le déterminant non
nul de \(\mathsf M_{\rm CKM}\) permet d’utiliser son inverse,
présenté dans \eqref{eq:ii-ckm-inverse}. La positivité de la métrique
découle de
\(x^*\mathsf M_{\rm CKM}^*\mathsf M_{\rm CKM}x
=\|\mathsf M_{\rm CKM}x\|^2>0\) pour \(x\ne0\).
\end{proof}

L’isométrie \(F\) conserve la métrique de comptage de
\(\boldsymbol\theta=\mathsf M_{\rm CKM}x_\theta\) ; son transport
aux coordonnées \(x_\theta\) donne la métrique du corollaire. Cette
distinction évite d’identifier une base de générateurs angulaires à une
base orthonormée de leurs paramètres.

\subsection{Certification exacte et composition unitaire finie}

Les racines positives du facteur polaire peuvent être vérifiées au moyen de leurs
projecteurs, en conservant l’arithmétique exacte. Nous définissons
\begin{equation}
 B=\frac16IP_\sigma A_R^*D,\qquad
 H=D^*GD=B^*B,\qquad L=H^{-1}B^*.
 \label{eq:ii-witt-radical-free}
\end{equation}
Les identités
\[
 LB=I_3,\qquad L(B\mathsf M_{\rm CKM})=\mathsf M_{\rm CKM},\qquad
 \mathsf M_{\rm CKM}^{-1}L(B\mathsf M_{\rm CKM})=I_3
\]
s’évaluent dans \(\mathbb Q(\sqrt5)\). Si \(E_D=D^*ED\),
\[
 H^{-1/2}=\lambda_s^{-1/2}E_D+
          \lambda_t^{-1/2}(I_3-E_D),\qquad F=BH^{-1/2}.
\]
Les produits \(E_D^2=E_D\), \(E_D(I_3-E_D)=0\) et la
décomposition de \(H\) prouvent cette formule sans approcher ses
radicaux. Les 132 lignes de \(B\) sont déterminées par
\eqref{eq:ii-witt-F-coordinates} et par la règle de construction des
hexades. Le certificat exécutable de la livraison conserve ces lignes,
celles de \(B\mathsf M_{\rm CKM}\), l’inverse à gauche \(L\),
les projecteurs et les étiquettes de chaque hexade. Ses 906 contrôles
exacts vérifient cette composition et ses identités d’orientation,
d’incidence et de récupération.

La conversion en radians s’effectue avant la réalisation trigonométrique :
\begin{equation}
 \boldsymbol\vartheta=
 \mathsf M_{\rm CKM}
 \begin{pmatrix}(\pi_{\rm HMT}/180)A_{\rm deg}\\
 (\pi_{\rm HMT}/180)C^*_{\rm deg}/6\\\Delta_4\end{pmatrix},
 \qquad
 \delta=\frac{\pi_{\rm HMT}}{180}\,9A_{\rm deg}.
 \label{eq:ii-witt-angular-units}
\end{equation}
La réalisation finie préserve l’ordre
\(V_{\rm CKM}=R_{23}U_{13}(\delta)R_{12}\) de
\eqref{eq:ii-ckm-factorization}. Pour préciser l’information supplémentaire
de la phase complexe, soit \(E_{ab}\) la matrice élémentaire et définissons
\[
 J_{13}=E_{13}-E_{31},\qquad S_{13}=E_{13}+E_{31},\qquad
 X_{13}(\delta)=\cos\delta\,J_{13}-i\sin\delta\,S_{13}.
\]
Alors \(X_{13}(\delta)^*=-X_{13}(\delta)\) et
\(U_{13}(\delta)=\exp(\vartheta_{13}X_{13}(\delta))\).
En effet, \(X_{13}(\delta)^2=-(E_{11}+E_{33})\), et son exponentielle
reproduit le bloc unitaire \(13\) de
\eqref{eq:ii-ckm-factorization}. Le transport extérieur conserve les
générateurs orientés ; la composante imaginaire symétrique conserve
l’information de phase. Le produit ordonné des trois facteurs
finis maintient ces deux contributions et leur ordre de composition.

\subsection{Résultat structurel et portée de la composition}

La construction fournit un transport isométrique équivariant depuis
les trois coordonnées angulaires jusqu’à l’incidence de Witt, ainsi qu’une
reconstruction explicite de \(A_{\rm deg}\), \(C^*_{\rm deg}/6\) et
\(\Delta_{\rm deg}\). Elle conserve l’orientation extérieure, le secteur
dodécaphasique, la face marquée du sceau et le changement conjoint de jauge.
En particulier, la réalisation d’incidence représente le même contenu
angulaire sur un autre support et permet de retrouver ses trois coordonnées.

La génération antérieure de \(K\), la sélection des fonctionnelles
sectorielles et l’identification physique des amplitudes conservent leurs
démonstrations propres. Les identités de récupération prouvées ici
établissent la fidélité de la composition ; leur utilisation ultérieure maintient
ces antécédents explicites. Le développement complémentaire des plans
quartiques de mémoire demeure une autre construction de l’état, et ne
s’interpose pas comme exigence supplémentaire de l’alignement réalisé dans
\eqref{eq:ii-witt-F}.
